\documentclass[border=0pt]{standalone}
% The shared preamble of every docs diagram (docs/diagrams/<name>.tex).
% dvisvgm's driver first, so patterns and opacity survive into the SVG.
\def\pgfsysdriver{pgfsys-dvisvgm.def}
\usepackage[siunitx,RPvoltages]{circuitikz}
\usepackage{xcolor}
% Latin Modern: outline fonts for every glyph a label uses, µ included.
\usepackage[T1]{fontenc}
\usepackage{lmodern}
\usetikzlibrary{arrows.meta,backgrounds,calc,patterns,positioning}
\tikzset{
  % A white card, so diagrams read in mdBook's dark themes too.
  every picture/.append style={show background rectangle,
    background rectangle/.style={fill=white,inner sep=4pt}},
  % Layout geometry: layers told apart by fill and hatching, never colour alone.
  layer a/.style={draw=black,line width=0.6pt,fill=gray!20},
  layer b/.style={draw=black,line width=0.6pt,fill=gray!45},
  layer c/.style={draw=black,line width=0.6pt,fill=white,postaction={pattern=north east lines}},
  layer d/.style={draw=black,line width=0.6pt,fill=white,postaction={pattern=crosshatch dots}},
  marker/.style={draw=black,dashed,line width=0.6pt},
  % A measured distance and its label.
  dim/.style={{Latex[length=2mm]}-{Latex[length=2mm]},line width=0.5pt},
  dim label/.style={font=\small,fill=white,inner sep=1pt},
  % Where a rule fails: thick, dashed and labelled, not only red.
  violation/.style={draw=red!75!black,line width=1.4pt,dash pattern=on 3pt off 2pt},
  violation label/.style={font=\small\bfseries,text=red!75!black},
  caption/.style={font=\small},
}

% One netlist line: node name #1 at height #2, text #3.
\newcommand{\nl}[3]{\node[anchor=west,font=\ttfamily\footnotesize,inner sep=1pt] (#1) at (0,#2) {#3};}
\begin{document}
\begin{circuitikz}[font=\small]
  \node[font=\bfseries,anchor=west] at (0,5.0) {top.v, the top level};
  \nl{v0}{4.4}{module top (A, Y, VDD, VSS);}
  \nl{v1}{3.95}{\ \ input A;}
  \nl{v2}{3.5}{\ \ output Y;}
  \nl{v3}{3.05}{\ \ inout VDD, VSS;}
  \nl{v4}{2.6}{\ \ inv X1 (.A(A), .Y(Y), .VDD(VDD), .VSS(VSS));}
  \nl{v5}{2.15}{endmodule}
  \node[font=\bfseries,anchor=west] at (0,1.45) {inv.spice, the cell library};
  \nl{s0}{0.85}{.subckt inv A Y VDD VSS}
  \nl{s1}{0.4}{M1 Y A VSS VSS nfet W=1u L=0.15u}
  \nl{s2}{-0.05}{M2 Y A VDD VDD pfet W=2u L=0.15u}
  \nl{s3}{-0.5}{.ends}

  \node[font=\bfseries] at (11.2,5.6) {Circuit it reads to};
  \node[pmos] (p) at (11.2,3.2) {};
  \node[right=2pt,align=left] at (p.east) {X1.M2 pfet\\ w \SI{2}{\micro\metre}};
  \draw (p.S) -- ++(0,0.3) node[vcc] {VDD};
  \node[nmos] (n) at (11.2,1.0) {};
  \node[right=2pt,align=left] at (n.east) {X1.M1 nfet\\ w \SI{1}{\micro\metre}};
  \draw (p.D) -- (n.D);
  \draw (p.D) ++(0,-0.35) -- ++(2.0,0) node[right] {Y};
  \draw (n.S) node[ground] {} node[right=6pt,yshift=-4pt] {VSS};
  \draw (p.G) -- ++(-0.6,0) coordinate (a) -- (a |- n.G) -- (n.G);
  \draw (a |- p.G) node[left] {A};
  \draw[marker] (10.4,-0.55) rectangle (13.3,4.1);
  \node[caption,anchor=south east] at (13.3,-0.55) {X1 (inv)};

  \nl{r0}{-1.5}{\ \ inv X1 (.A(A), .Y(Y), .VDD(VDD), .VSS(1\textquotesingle b0));}
  \draw[violation] (-0.15,-1.85) rectangle (8.1,-1.15);
  \node[violation label,anchor=north west] at (-0.15,-2.0) {refused: the constant 1\textquotesingle b0 in a connection (line 5)};
\end{circuitikz}
\end{document}
